\section{Acquired continuation geometry and causal transfer}\label{sec:transfer}
We first specify the lower certificate, which depends only on the physical experiment and its fixed continuation, not on a proposed encoder.

\begin{definition}[Dominated continuation cut]\label{def:cut}
At a deterministic cut $s\le T$, let $H$ be the actual prefix and $U$ the complete physical pre-decision suffix, including all post-cut public actions, phases, costs and reports. The complete decision history is generated by $(H,U)$. The fixed exploration supplies their joint law $\Lambda(dh,du)$. A dominated continuation certificate consists of a probability measure $\eta$ on the suffix space and a finite measure $\underline\mu$ on prefix histories such that
\begin{equation}\label{eq:domination}
 \underline\mu(dh)\eta(du)\le\Lambda(dh,du).
\end{equation}
Choose a bounded jointly Borel version $z(h,u)=\E[Y\mid H=h,U=u]$ and set
\[
 Z_s(h)=z(h,\cdot)\in L^2(\eta;\R^q),\qquad
 \xi_s=(Z_s)_*\underline\mu.
\]
The task norm supplies the inner product in this function space. All post-cut information supplied to a decoder is contained in $U$ and independent algorithm coins; no earlier tape or additional data-dependent public register remains outside its counted label.
\end{definition}

The product inequality is a raw-law minorization, not independence of the actual suffix. For example a density $j(u\mid h)\ge c$ relative to $\eta$ on an actual prefix event $E$ gives $\underline\mu=c\,\law(H)|_E$. An independent suffix has $c=1$. In the absence of such a nonzero certificate the corresponding bound is zero, not an assertion that the experiment is easy. Domination ensures that versions of $z$ give the same element for $\underline\mu$-almost every $h$. Standard Borel spaces have countably generated sigma fields modulo a fixed probability, so $L^2(\eta)$ is separable. Approximation of a bounded jointly measurable function by simple rectangle functions, followed by dominated convergence, makes $h\mapsto Z_s(h)$ a Borel map into this Hilbert space.

\begin{lemma}[Projection and the continuation cut]\label{lem:cut}
For the fixed task in \eqref{eq:task},
\begin{equation}\label{eq:checkpoint}
 R_\cp(T,M)=B_T+q_M(\nu_T).
\end{equation}
Every $M$-label online predictor and every certificate in Definition~\ref{def:cut} satisfy
\begin{equation}\label{eq:cutlower}
 R_\on(T,M)-B_T\ge q_M(\xi_s).
\end{equation}
Both assertions allow independent shared and private randomization. They use the same terminal task and baseline.
\end{lemma}
\begin{proof}
Conditional projection gives $\E\norm{Y-A}^2=B_T+\E\norm{P_T-A}^2$. Fix all algorithm coins, which are independent of the fixed exploration. A checkpoint readout has at most $M$ centers. Randomized readouts may alternatively be replaced by conditional means, and nearest-center selection is Borel. Taking approximating center lists proves \eqref{eq:checkpoint}; averaging independent seeds does not change $\nu_T$.

At cut $s$ write its retained label as $i(H)$. Give the suffix decoder all of $U$, even when its original online implementation would retain less. Conditional on fixed coins, including all future private decoder coins, composing the finite sequence of later Borel updates gives an output $a_{i(H)}(U)$ for at most $M$ Borel response functions. Here $U$ contains all post-cut public actions, phases, costs and reports, not observation values alone. Projection and \eqref{eq:domination} give
\[
 R-B_T\ge\int\norm{z(h,u)-a_{i(h)}(u)}^2
                   \underline\mu(dh)\eta(du)
 \ge\int\min_i\norm{Z_s(h)-a_i}_{L^2(\eta)}^2\underline\mu(dh).
\]
Outputs can be projected into a fixed bounded convex set containing the conditional means, so all center functions belong to $L^2(\eta)$. The last expression is at least $q_M(\xi_s)$. Averaging the coins proves the assertion. No different continuation task is inserted into the loss, and no product structure of $\Lambda$ beyond the lower certificate is used.
\end{proof}

We next state a constructive upper certificate. It is convenient to allow a finite declared phase as in Section~\ref{sec:problem}; otherwise take one fiber. All constants below are uniform over phases and the specified strategies. A code has at most $M$ representatives in each phase. A cut lower must replace $M$ by the size of the full accessible tuple when its phase itself carries data.

Let $(S_v,d_v)$ be separable metric predictive fibers, with Borel envelope $w\ge1$. The update along a common report and physical action is $F_t$. A Borel relation $\mathcal A_v$ satisfies $x\in\mathcal A_v(x)$ and $w(z)\le w(x)$ for $z\in\mathcal A_v(x)$. It is fixed independently of the budget. Finite Borel covers have representatives $c_i^v$ and first-cell selectors $Q_M^v$ such that
\begin{equation}\label{eq:cover}
 Q_M^v(x)\in\mathcal A_v(x),\qquad
 d_v(x,Q_M^v(x))\le r_M w(x).
\end{equation}
For $s=km$, write $\mathscr F_{s,j}(z)$ for $j$ exact updates along the next physical word and the actions selected from that physical transcript. For every possible block-start past, true state $x$ and candidate $z$ in its named fiber, require one-step pathwise Lipschitz gain at most $L\ge1$ and
\begin{equation}\label{eq:pair}
 \E_x[d(\mathscr F_{s,m}(x),\mathscr F_{s,m}(z))^2\mid\mathcal H_s]
 \le\kappa d(x,z)^2,\qquad 0\le\kappa<1.
\end{equation}
The notation $d$ denotes the appropriate common fiber metric. For \emph{every} adapted sequence $z_0=z$, $v_j=F_{s+j}(z_{j-1},a_{s+j},Y_{s+j})$, $z_j\in\mathcal A(v_j)$, require
\begin{align}
 \E_x[w(z_m)^2\mid\mathcal H_s]&\le a w(z)^2+b,\label{eq:drift}\\
 \E_x[w(v_j)^2\mid\mathcal H_s]&\le A w(z)^2+B\quad(1\le j\le m),\label{eq:within}
\end{align}
where $0\le a<1$ and $b,A,B\ge0$. The conditional inequalities continue to hold when past independent algorithm coins are included. They quantify over arbitrary relation-preserving adapted perturbations, rather than assuming the performance of a preconstructed quantizer. Numerical implementations must still choose the same codebook and relation, with local error at most $r_Mw(v_j)+u$.

\begin{theorem}[Causal acquired-geometry and continuation transfer]\label{thm:transfer}
Fix a prepared experiment, a deterministic horizon and encoder-independent exploration. Suppose \eqref{eq:cover}--\eqref{eq:within} and a legal seed with $\E w(\widehat S_0)^2\le H_0$ and $\norm{d(S_0,\widehat S_0)}_2\le e_0$. Assume $P_T=f_T(S_T)$ with Lipschitz constant $L_f$ in the task norm. Put
\begin{align*}
 H&=\max\{H_0,b/(1-a)\},& H_*&=AH+B,\\
 \Lambda_j&=\sum_{i=0}^{j-1}L^i,& \Delta_M&=r_M\sqrt{H_*}+u,\\
 E_*&=\max\{e_0,\Lambda_m\Delta_M/(1-\sqrt\kappa)\},&
 E_M&=L^{m-1}E_*+\Lambda_{m-1}\Delta_M.
\end{align*}
For a one-fiber $M$-label model, or a model in which $M$ below counts the whole accessible tuple, define
\[
 W_M=\max\{q_M(\nu_T),\ \sup_{s,\,\text{certificates}}q_M(\xi_s)\}.
\]
All certificate measures are unnormalized: scaling $\underline\mu$ scales $q_M(\xi_s)$ by the same factor.
There is a causal predictor quantized after every raw call, with terminal readout error at most $v$ in root mean square, such that
\begin{equation}\label{eq:central}
 \begin{gathered}
 R_\cp(T,M)=B_T+q_M(\nu_T),\\
 B_T+W_M\le R_\on(T,M)
       \le B_T+(L_fE_M+v)^2.
 \end{gathered}
\end{equation}
For phase-indexed covers the upper uses the stated $M$ labels per phase and their charged full tuple; the lower always uses that tuple's cardinality.

The lower can also include every conditional continuation width of Theorem~\ref{thm:cutduality}. On the sequential-chart class of Theorem~\ref{thm:chart}, their supremum and the optimal recursively achievable excess are comparable; outside that class no sufficiency is claimed.

For any actual measure $\xi$, either $\nu_T$ or a dominated continuation measure, a submeasure of mass $\zeta>0$ with
\begin{equation}\label{eq:smallball}
 \sup_c\underline\xi(B(c,s_M))\le\zeta/(2M)
\end{equation}
contributes at least $\zeta s_M^2/2$ to the lower. Thus a matching actual terminal witness with $s_M\asymp r_M$, mass bounded below, $e_0,u,v=O(r_M)$ and fixed certificate constants gives matching checkpoint and online excess of order $r_M^2$. A larger continuation width instead certifies a separation from the terminal checkpoint curve. All terms refer to the one fixed task and actual preparation.
\end{theorem}
\begin{proof}
The exact selector in \eqref{eq:cover} is Borel and yields the recursion $\widehat S_t=Q_M(F_t(\widehat S_{t-1},a_t,Y_t))$. The numerical routine obeys the same relation. Equation~\eqref{eq:drift} proves inductively that the envelope second moment is at most $H$ at block endpoints; \eqref{eq:within} bounds each pre-rounding moment by $H_*$. Hence each local insertion of error has $L^2$ norm at most $\Delta_M$.

Start an exact shadow at $\widehat S_s$ for a block and feed it the same physical word. Telescoping the local insertions, with all subsequent pathwise gains included, gives
\[
 d(\mathscr F_{s,m}(\widehat S_s),\widehat S_{s+m})
 \le\sum_{j=1}^m L^{m-j}d(v_j,z_j).
\]
Conditional use of \eqref{eq:pair}, then Minkowski, shows that the root errors satisfy
\[
 e_{s+m}\le\sqrt\kappa e_s+\Lambda_m\Delta_M,
 \qquad e_{s+j}\le L^je_s+\Lambda_j\Delta_M\quad(0\le j<m).
\]
The first preserves $e_{km}\le E_*$ and the second gives $e_T\le E_M$. The shadow is a proof device: no block of reports, real posterior or second state is retained by the machine. Lipschitz readout, root-error addition and conditional projection prove the upper.

Lemma~\ref{lem:cut} gives all the lower bounds simultaneously for this same experiment. Their supremum remains a lower bound. Finally $M$ open balls satisfying \eqref{eq:smallball} cover at most half the submeasure. Outside their union the squared distance from every list is at least $s_M^2$, proving the stated contribution. This is valid in the continuation Hilbert space as well as in a finite-dimensional score space. The matching assertions follow directly; no ambient-dimensional volume has been assumed.
\end{proof}

\begin{corollary}[A necessary compatibility condition]\label{cor:necessary}
For a sequence of fixed-task experiments and budgets, an asserted bound
$R_\on(T,CM)-B_T\le C' q_M(\nu_T)$ implies
$W_{CM}\le C' q_M(\nu_T)$, with the same specified full-state cut. A family for which this inequality fails cannot have that claimed online/checkpoint matching, even if its terminal marginal has small covering number. Here $CM$ already denotes the entire data-dependent tuple at that cut. A separately stored phase of $Q$ values must be accounted for as $CMQ$ unless its value is deterministic at that cut; the constant $C$ cannot silently absorb a growing phase.
\end{corollary}
\begin{proof}
Apply the lower side of \eqref{eq:central}, or Lemma~\ref{lem:cut} alone when no stability certificate is available.
\end{proof}

\begin{example}[A quantitative delayed-query separation]\label{ex:query}
Prepare $d$ independent fair bits $X_1,\ldots,X_d$ and report them in $d$ paid calls. After the last retained cut report an independent uniform index $J\in\{1,\ldots,d\}$. The task is to predict $Y=X_J$ before its terminal audit. A checkpoint reading the full history after $J$ needs two labels for zero risk. If an online predictor had at most $M$ labels immediately before $J$, then
\begin{equation}\label{eq:randomaccess}
 R_\on\ge\tfrac12 h_2^{-1}\bigl((1-\log_2 M/d)_+\bigr),
\end{equation}
where $h_2^{-1}$ is the inverse binary entropy on $[0,1/2]$. In particular retaining at most a fixed fraction of $d$ bits gives a positive error bounded away from zero. Exact prediction requires $M\ge2^d$ and storing the bit string attains it.
\end{example}
\begin{proof}
The continuation measure is the law of $j\mapsto X_j$ in the uniform $L^2$ norm. For an independent shared seed $R$ and retained label $C$, $I(X;C\mid R)\le\log_2 M$. Independence of the bits and conditional subadditivity give
$\sum_j H(X_j\mid C,R)\ge d-\log_2M$.
Let $e$ be the average posterior bit-classification error. The two inequalities used are
\[
 h_2(e)\ge(1-\log_2M/d)_+,\qquad
 p(1-p)\ge p/2\quad(0\le p\le1/2).
\]
The first is concavity of binary entropy; the second bounds Bernoulli posterior variance by its classification error. Conditional projection proves \eqref{eq:randomaccess}. Zero error forces distinct positive-probability labels for all strings. This is the classical entropy mechanism behind random-access bounds~\cite{nayak}, used here to distinguish two information cuts of a single prediction task, not claimed as a new coding inequality.
\end{proof}

\begin{proposition}[Flagged violations of an admissibility relation]\label{prop:repair}
Suppose a finite implementation agrees with a routine satisfying the upper hypotheses until a flagged first failure, and, conditional on the current past and on no earlier flagged failure, its probability of failing at step $t$ is at most $\beta_t$. A repair routine using the same codebook can always be selected on a failure to satisfy the original relation and local-error bound. For a loss in $[0,L_*]$, the actual implementation has risk at most the repaired upper plus $L_*\min(1,\sum_{t\le T}\beta_t)$.
\end{proposition}
\begin{proof}
Drive the repaired routine by the same physical reports and coins; at a failure choose the certified selector instead. Encoder-independent exploration makes the physical law unchanged. The two retained trajectories agree until the first failure. A first-failure union bound gives its probability at most $\sum_t\beta_t$, and bounded loss gives the assertion. The reference estimate is unconditional; it is not incorrectly conditioned on the absence of failure. The proposition weakens implementation accuracy to a probabilistic certificate, but does not assert the existence of global relation-preserving covers from local data alone.
\end{proof}
